% --- GRPO and GSPO ---
% --- Source: Lecture-11-12\_NLP\_Recent\_Advancements (sections/recent-advance/deepseek-grpo.tex); DeepSeekMath / DeepSeek-R1 ---
\begin{frame}{GRPO Overview}
    \begin{itemize}
        \item \textbf{GRPO}: Group Relative Policy Optimization, an evolution of PPO that is more stable, efficient, and converges faster
        \item Combines reward and preference signals, using single model for both
        \item \textcolor{blue}{Improves safety and alignment of open-weight models}
        \item Introduced in DeepSeekMath, used in DeepSeek-MoE and other models
    \end{itemize}
\end{frame}

\begin{frame}[allowframebreaks]{GRPO Objective}
    \begin{itemize}
        \item GRPO avoids the need for a separate critic model by estimating the baseline from group scores.
        \item For each question $q$, sample a group of outputs $\{o_1, o_2, \ldots, o_G\}$ from the old policy $\pi_{\theta_{\text{old}}}$.
        \item The policy model $\pi_\theta$ is optimized by maximizing:
        \begin{equation}
        \begin{split}
            J_{\text{GRPO}}(\theta) =\; & \mathbb{E}_{q \sim P(Q),\, \{o_i\}_{i=1}^G \sim \pi_{\theta_{\text{old}}}(O|q)} \Bigg[
            \frac{1}{G} \sum_{i=1}^G \Bigg(
                \min \Bigg(
                    \frac{\pi_\theta(o_i|q)}{\pi_{\theta_{\text{old}}}(o_i|q)} A_i, \\
                & \qquad\qquad
                    \text{clip}\left(\frac{\pi_\theta(o_i|q)}{\pi_{\theta_{\text{old}}}(o_i|q)}, 1-\epsilon, 1+\epsilon\right) A_i
                \Bigg)
                - \beta D_{\text{KL}}(\pi_\theta || \pi_{\text{ref}})
            \Bigg)
            \Bigg]
        \end{split}
        \end{equation}
    \framebreak
        \item KL divergence term:
        \begin{equation}
            D_{\text{KL}}(\pi_\theta || \pi_{\text{ref}}) =
            \frac{\pi_{\text{ref}}(o_i|q)}{\pi_\theta(o_i|q)} - \log \frac{\pi_{\text{ref}}(o_i|q)}{\pi_\theta(o_i|q)} - 1
        \end{equation}
        \item Advantage $A_i$ is computed using group rewards $\{r_1, r_2, \ldots, r_G\}$:
        \begin{equation}
            A_i = \frac{r_i - \text{mean}(\{r_1, r_2, \ldots, r_G\})}{\text{std}(\{r_1, r_2, \ldots, r_G\})}
        \end{equation}
        \item $\epsilon$ and $\beta$ are hyper-parameters.
    \end{itemize}
\end{frame}

\begin{frame}{Why GRPO Is Popular for Open Reasoning Models}
  \begin{itemize}\itemsep0.4em
    \item[\textcolor{red}{$\blacktriangleright$}] \textbf{Fewer models in memory:} policy + reference (no critic, no RM when rules suffice).
    \item[\textcolor{red}{$\blacktriangleright$}] \textbf{Often faster to train} than PPO-style RLHF on comparable hardware.
    \item[\textcolor{red}{$\blacktriangleright$}] Powers DeepSeekMath, DeepSeek-R1-Zero, and many open reproductions.
    \item[\textcolor{red}{$\blacktriangleright$}] With QLoRA / optimized kernels, 8B-scale GRPO fits consumer GPUs (order-of-magnitude below naive full fine-tune).
  \end{itemize}
  \vspace{0.35em}
  \small
  \begin{tabular}{@{}lccc@{}}
    \toprule
    \textbf{Setup (Llama 8B, indicative)} & \textbf{Naive} & \textbf{Optimized} & \textbf{Save} \\
    \midrule
    GRPO (full) & $\sim$510\,GB & $\sim$160\,GB & $\sim$68\% \\
    GRPO + QLoRA & --- & $<$10\,GB & $>$90\% \\
    \bottomrule
  \end{tabular}
\end{frame}

\begin{frame}{GRPO vs.\ PPO (Takeaways)}
  \small
  \begin{tabular}{@{}p{0.22\textwidth}p{0.36\textwidth}p{0.36\textwidth}@{}}
    \toprule
    & \textbf{PPO (RLHF)} & \textbf{GRPO (RLVR)} \\
    \midrule
    Baseline & Critic / value net & Mean reward in a group \\
    Reward & Learned RM & Rules / verifiers \\
    Models in memory & Policy, ref., RM, critic & Policy, ref. (typical) \\
    Best for & Chat alignment & Math, code, logic \\
    \bottomrule
  \end{tabular}
  \vspace{0.5em}
  \begin{itemize}\itemsep0.3em
    \item[\textcolor{red}{$\blacktriangleright$}] GRPO does \textbf{not} replace DPO---DPO still wins when you only have preference pairs and no verifier.
  \end{itemize}
\end{frame}

\begin{frame}{Motivation: Why GRPO Breaks at Scale}
    \begin{itemize}\itemsep0.35em
      \item[\textcolor{red}{$\blacktriangleright$}] Long RL runs on huge models need \textbf{stable} updates---unstable RL $\Rightarrow$ \textbf{model collapse} (irreversible).
      \item[\textcolor{red}{$\blacktriangleright$}] \textbf{GRPO} uses \textbf{token-level} importance ratios and \textbf{per-token} clipping on long reasoning chains.
      \item[\textcolor{red}{$\blacktriangleright$}] Clipping can still leave many tokens with skewed weights; effects compound over training.
      \item[\textcolor{red}{$\blacktriangleright$}] On \textbf{MoE} models, token-level updates amplify \textbf{routing churn} (experts change mid-training)---often needs extra fixes (e.g.\ Routing Replay).
    \end{itemize}
\end{frame}

% --- Source: Zheng et al.\ (2025) arXiv:2507.18071; Qwen blog https://qwen.ai/blog?id=gspo ---

\begin{frame}{GSPO: Group Sequence Policy Optimization}
  \begin{itemize}\itemsep0.4em
    \item[\textcolor{red}{$\blacktriangleright$}] Rewards are usually \textbf{sequence-level} (correct answer, tests passed)---optimization should match that granularity.
    \item[\textcolor{red}{$\blacktriangleright$}] Same group-relative advantages as GRPO: compare $G$ completions on the same prompt.
    \item[\textcolor{red}{$\blacktriangleright$}] Design goals (\href{https://qwen.ai/blog?id=gspo}{Qwen blog}):
    \begin{itemize}\itemsep0.2em
      \item \textbf{Stable}---especially for MoE RL;
      \item \textbf{Efficient}---better compute-to-gain than GRPO;
      \item \textbf{Infrastructure-friendly}---more robust to numeric precision issues in distributed RL.
    \end{itemize}
  \end{itemize}
\end{frame}

\begin{frame}{GSPO: Algorithm Architecture (Training Loop)}
  \small
  \begin{enumerate}\itemsep0.35em
    \item[\textcolor{red}{$\blacktriangleright$}] \textbf{Sample} a batch of prompts $x \sim \mathcal{D}$.
    \item[\textcolor{red}{$\blacktriangleright$}] For each $x$, draw a \textbf{group} $\{y_i\}_{i=1}^G$ from the old policy $\pi_{\theta_{\mathrm{old}}}$.
    \item[\textcolor{red}{$\blacktriangleright$}] \textbf{Reward} each full response (verifier / rule checker) $\Rightarrow$ group-relative advantage $\hat{A}_i$ (same formula as GRPO).
    \item[\textcolor{red}{$\blacktriangleright$}] Build a \textbf{sequence-level} importance weight $s_i(\theta)$ from $\pi_\theta(y_i|x)$ vs.\ $\pi_{\theta_{\mathrm{old}}}(y_i|x)$.
    \item[\textcolor{red}{$\blacktriangleright$}] \textbf{Clip and optimize once per sequence} (not token-by-token); optional KL to $\pi_{\mathrm{ref}}$ as in PPO/GRPO stacks.
    \item[\textcolor{red}{$\blacktriangleright$}] \textbf{Variant:} GSPO-token allows per-token advantages when needed (e.g.\ multi-turn RL).
  \end{enumerate}
\end{frame}

\begin{frame}{GRPO vs.\ GSPO}
  \small
  \begin{tabular}{@{}p{0.26\textwidth}p{0.34\textwidth}p{0.34\textwidth}@{}}
    \toprule
    & \textbf{GRPO} & \textbf{GSPO} \\
    \midrule
    Importance ratio & Per-token & Sequence-level $s_i(\theta)$ \\
    Clipping & Token-by-token & Once per sequence \\
    Advantage & Group-relative & Group-relative (same) \\
    Reward alignment & Mismatch on long CoT & Matches sequence rewards \\
    MoE RL & Often unstable & Stable in Qwen3-scale runs \\
    Typical use & DeepSeek-R1, dense RL & Qwen3, large MoE RL \\
    \bottomrule
  \end{tabular}
  \vspace{0.45em}
  \begin{itemize}\itemsep0.3em
    \item[\textcolor{red}{$\blacktriangleright$}] GSPO can clip \textbf{more tokens} in aggregate yet train \textbf{more efficiently}---lower-variance signal.
  \end{itemize}
\end{frame}
